\documentclass[11pt,a4paper]{article}
\usepackage[margin=1in]{geometry}
\usepackage{amsmath,amssymb,amsthm}
\usepackage{graphicx}
\usepackage{microtype}
\usepackage[hidelinks]{hyperref}

\newtheorem{theorem}{Theorem}
\newtheorem{corollary}[theorem]{Corollary}
\newcommand{\T}{\mathbb{T}}
\newcommand{\C}{\mathbb{C}}
\newcommand{\Z}{\mathbb{Z}}
\newcommand{\E}{\mathcal{E}}

\title{Positive non-atomic minimal extrapolations are never unique}
\author{Full answer to the explicit uniqueness questions in arXiv:1601.05761}
\date{17 August 2026}

\begin{document}
\maketitle

\noindent\textbf{Status.}
Candidate full solution, likely valid, pending human review.  The two explicit
uniqueness questions in Examples~3.8 and~3.9 of the source are both answered
negatively.  The general theorem also settles every positive non-atomic case
when the sampled frequency set is finite and contains zero.

\section{Source questions}

Benedetto and Li minimize total variation among measures having prescribed
Fourier coefficients on a finite set $\Lambda\subset\Z^d$~\cite{BenedettoLi}.
On source PDF page~20, Example~3.8 says that their theory cannot determine
whether a positive continuous singular Cantor measure $\sigma_q$ is the unique
minimal extrapolation.  Example~3.9 asks the same question for
$\mu=\sigma_A+\sigma_B$, where $\sigma_A$ and $\sigma_B$ are Haar surface
measure on the horizontal circles $x_2=0$ and $x_2=1/2$ in $\T^2$.

\begin{center}
\includegraphics[width=\textwidth]{figures/open_problem_crop.png}
\end{center}

There is a minor internal inconsistency in Example~3.9: the printed set is
$\{-2,-1,\ldots,1\}^2$, while the displayed set $\Gamma$ contains frequencies
with second coordinate $2$.  The explicit answer below matches the data on the
larger rectangle $\{-2,-1,0,1,2\}^2$, and therefore covers both readings.

\section{Proof intuition}

The zero Fourier coefficient is decisive.  For a positive measure $\mu$, it
equals $\|\mu\|$, so every extrapolation has norm at least $\|\mu\|$.  Thus
every positive measure sharing the sampled moments is automatically another
minimal extrapolation.  Finitely many Fourier constraints impose only
finitely many real linear equations on a bounded density perturbation $h$.
A non-atomic measure can be split into more positive-measure pieces than there
are equations, leaving a nonzero kernel vector.  After scaling $h$, both
$(1+h)\mu$ and $(1-h)\mu$ remain positive and share all sampled Fourier data.

For the surface-measure example this abstract perturbation can be made fully
explicit: replace uniform measure in the $x_1$ direction by the uniform
five-point rule on the fifth roots of unity.  It integrates every Fourier mode
$m_1=-2,-1,0,1,2$ exactly.

\section{General non-uniqueness theorem}

For a finite positive measure $\mu$ on $\T^d$, write
\[
 \widehat\mu(m)=\int_{\T^d}e^{-2\pi i m\cdot x}\,d\mu(x).
\]

\begin{theorem}\label{thm:general}
Let $\mu$ be a nonzero positive non-atomic finite measure on $\T^d$, and let
$\Lambda\subset\Z^d$ be finite with $0\in\Lambda$.  Then $\mu$ is a minimal
extrapolation of $\widehat\mu|_\Lambda$, but it is not the unique minimal
extrapolation.  In fact, there are distinct positive measures $\mu_+$ and
$\mu_-$ such that
\[
 \widehat{\mu_+}(m)=\widehat{\mu_-}(m)=\widehat\mu(m)
 \quad(m\in\Lambda),
 \qquad
 \|\mu_+\|=\|\mu_-\|=\|\mu\|.
\]
\end{theorem}

\begin{proof}
Let $r=2|\Lambda|+1$.  Non-atomicity permits a measurable partition
$\T^d=E_1\sqcup\cdots\sqcup E_r$ modulo $\mu$-null sets, with
$\mu(E_j)>0$ for every $j$.  For $c=(c_1,\ldots,c_r)\in\mathbb{R}^r$, put
$h_c=\sum_jc_j\mathbf{1}_{E_j}$.  The conditions
\begin{equation}\label{eq:constraints}
 \int_{\T^d} e^{-2\pi i m\cdot x}h_c(x)\,d\mu(x)=0
 \qquad(m\in\Lambda)
\end{equation}
form at most $2|\Lambda|$ homogeneous real linear equations in the $r$
real unknowns $c_j$.  Hence they admit a nonzero solution $c$.  Rescale it so
that $\|h_c\|_\infty\le 1/2$, and write $h=h_c$.

Define $d\mu_+=(1+h)d\mu$ and $d\mu_-=(1-h)d\mu$.  These measures are
positive and distinct, and~\eqref{eq:constraints} gives their asserted Fourier
data.  Since $0\in\Lambda$, every extrapolation $\nu$ satisfies
\[
 \|\nu\|\ge |\widehat\nu(0)|
   =|\widehat\mu(0)|=\mu(\T^d)=\|\mu\|.
\]
On the other hand, positivity and the zero-frequency equality give
$\|\mu_\pm\|=\mu_\pm(\T^d)=\mu(\T^d)$.  Thus $\mu$, $\mu_+$, and $\mu_-$
are all minimal extrapolations, proving non-uniqueness.
\end{proof}

\begin{corollary}[Answer to Example 3.8]
Under the standing finite-frequency assumption of the source, its positive
continuous singular Cantor measure $\sigma_q$ is not the unique minimal
extrapolation from any finite $\Lambda$ containing $0$.
\end{corollary}

\begin{proof}
The source constructs $\sigma_q$ as a positive continuous measure of mass one.
It is therefore non-atomic, and Theorem~\ref{thm:general} applies.
\end{proof}

\section{Explicit answer to the surface-measure example}

Let $\sigma_A$ and $\sigma_B$ be normalized Haar measure on
$A=\{x_2=0\}$ and $B=\{x_2=1/2\}$, and set $\mu=\sigma_A+\sigma_B$ as in
Example~3.9.  Define the positive atomic measure
\[
 \nu=\frac15\sum_{j=0}^{4}
 \left(\delta_{(j/5,0)}+\delta_{(j/5,1/2)}\right).
\]

\begin{corollary}[Answer to Example 3.9]
The source measure $\mu$ is not the unique minimal extrapolation.  The measure
$\nu$ is a distinct minimal extrapolation for the printed frequency set and
also for the symmetric set $\{-2,-1,0,1,2\}^2$ suggested by the displayed
$\Gamma$.
\end{corollary}

\begin{proof}
For $m=(m_1,m_2)\in\Z^2$,
\[
 \widehat\mu(m)=\mathbf{1}_{\{m_1=0\}}
                    \bigl(1+(-1)^{m_2}\bigr),
\]
whereas
\[
 \widehat\nu(m)=
 \left(\frac15\sum_{j=0}^{4}e^{-2\pi i m_1j/5}\right)
 \bigl(1+(-1)^{m_2}\bigr).
\]
The parenthesized roots-of-unity sum is one when $5$ divides $m_1$ and zero
otherwise.  Hence $\widehat\nu(m)=\widehat\mu(m)$ whenever $|m_1|\le2$,
with no restriction on $m_2$.  In particular the equality holds on either
reading of the source frequency set.

Both measures are positive with total mass two, while every extrapolation has
norm at least the absolute value $2$ of its zero Fourier coefficient.  Thus
both are minimal.  Finally, $\nu$ is atomic and $\mu$ is non-atomic, so they
are distinct.
\end{proof}

\section{Verification, novelty, and scope}

The accompanying script
\texttt{code/verify\_five\_point\_extrapolation.py} evaluates the explicit
Fourier formulas on $m_1=-2,\ldots,2$ and a larger range of $m_2$; the exact
roots-of-unity proof above is independent of the computation.

A bounded novelty search covered the run indexes, exact phrases from
Examples~3.8--3.9, the promised title ``Super-resolution of singular continuous
measures on $\T^d$'', and close arXiv searches for positive non-atomic minimal
extrapolations and finite Fourier moments.  The promised sequel remains listed
as ``in preparation'' in the later journal version and a 2018 dissertation;
no primary source giving these answers was found.  The underlying
finite-moment argument is elementary, so novelty confidence is moderate even
though correctness confidence is high.

The result fully settles the two explicit source uniqueness questions.  The
general theorem requires positivity, non-atomicity, a finite frequency set,
and inclusion of zero; it does not classify signed or complex singular
continuous measures.  Human review should verify the real dimension count in
\eqref{eq:constraints} and the scope attached to the word ``full''.

\begin{thebibliography}{1}
\bibitem{BenedettoLi}
J.~J. Benedetto and W.~Li,
\emph{Super-resolution by means of Beurling minimal extrapolation},
arXiv:1601.05761; Appl. Comput. Harmon. Anal. 48 (2020), 117--148.
\end{thebibliography}

\end{document}
